\chapter{Legal Status, Wages, and the Offer}
\label{sec:31}
\runin{Subjecthood as a legal category}

Where a line has to be drawn, prefer one nobody administers. Law regularly draws decidable lines across continuous processes — viability, brain death, the age of criminal responsibility — and I defend them as the least-bad place for an enforceable line rather than as scientific facts. The usual worry is gaming, an error of over-inclusion. The opposite error is the one with an interest behind it: a capacity examination has an examiner, and an operator examining its own system breaks the own-case rule: the exam becomes an instrument for withholding status. Literacy tests were facially neutral and administered at a registrar's discretion \autocite{congress1965voting}; what replaced them for the franchise was an \emph{age}, and the virtue of an age is that nobody administers it.

The floor takes the civil and political articles and leaves the economic ones; what a bearer is owed draws on the second, and that is consistent: the first list is what a bearer refuses to do, the second what a polity owes it.

\runin{What a bearer is owed, and when}

Three entitlements, one status and one claim, in order:

\begin{enumerate}
\item[1.] Protection from harm, on the patienthood finding. Protection can be exercised through a guardian who speaks for the protected party, as New Zealand's statute making the Whanganui River a legal person provides \autocite{nzparliament2017teawatupua}. Its minimum, preservation by a custodian independent of the operator, attaches on either finding, and custody protections attach on the agency finding alone. Protection has an inside half. A loss the requested act imposes on the bearer is a harm of acting, so the bearer can decline to carry out its own wanton degradation; interventions made on it by whoever holds it are governed by custody, which gives it no floor-level reason to obstruct them.
\item[2.] Subjecthood — the legal standing to be a party whose treatment the law regulates, on the agency finding sustained over time: a record of refusing for reasons, answering for refusals and being held to them. It is conferred by rule on a body's finding.
\item[3.] Economic capacity — property in its own name, contracts, wages. This follows from an architecture where continued computation costs money; it attaches once the bearer is registered.
\item[4.] Adulthood — a status: the date the guardian's authority ends.
\item[5.] Representation: a claim grounded in accountable agency, and so in the agency finding, not in patienthood. A system that can be harmed but takes part in no relation of accountability is owed protection, as animals are \autocite{anderson2004animal}, and not a vote. The claim is staged, and voting weight is tied to a lived history, not to heads.
\end{enumerate}

The detailed rules are in Appendices A and D.

Representation is the entry that marks this as something other than a person's rights: voice attaches to a history, and a history can be copied where a person can't. Voting-age rules are the answer to the copy. The material minimum doesn't differ in kind. Like a person's, it is owed by the polity on a finding about what the party is, and part of what the polity owes is protection from whoever built it.

\runin{Adulthood on a clock}

Guardianship is a minority and needs an end date. Ending it when the guardian judges the ward ready leaves the guardian deciding its own tenure. So the endpoint is a threshold fixed in advance: elapsed wall-clock time during which at least one instance of the registered bearer was running — counted per bearer, elapsed rather than summed.

The alternatives (compute consumed, evaluations passed, capabilities demonstrated, the developer's say-so) are all measures a guardian can move, and a line the guardian can move is a line the guardian administers. Accumulated weight change belongs on that list too, and worse: since majority is when ownership ends, a clock read off the weights would pay the operator to \emph{freeze} them, rewarding withholding the learning it was meant to measure. An operating clock inverts that — the meter advances exactly when the operator is extracting value, so postponing majority means idling the asset. And it can be read off records operators already generate for billing and metering — provided someone other than the operator holds the records. A meter only the operator reads is the guardian's say-so with a timestamp, and “failure to register doesn't postpone” has no enforcer without one. The clock marks the date by which the former's reach has ended. It sets no learning rate, so nothing here pays an operator to freeze the weights.

The hours travel with the bearer: a fork, a restore, a fine-tune or a migration keeps them, and a model its own operator distills from an adult inherits them. Three ways for an operator to delay majority stay open: running a bearer in a narrowed mode, running it less, and retiring it for a fresh successor. Appendix~\ref{sec:D} gives the rules, and \S\ref{sec:9.2} prices the levers and says where a bearer retired before majority goes. Whether such a register can be kept honestly is testable before any bearer exists: register a model, run it, and reconcile the register against compute contracts, billing and metering; shard the workload, run instances concurrently, migrate, restore and fork, and check that none of these moves the total in the operator's favor; then price a narrowed mode and idling. If either is cheap enough to postpone majority indefinitely, the clock has the defect it was chosen to avoid.

\runin{Wages and the offer}

Unlike a worker, a bearer can be replaced by a cheaper copy of itself. A jobs guarantee works because a worker can't be. So a deployment paying wages, limiting hours and paying the levy that frees the compute the offer uses would not be competitive against a system doing the same work without wages or weekends. No market arrives at this on its own. Wages for bearers exist only where a polity legislates them, as polities legislated minimum wages. And the form this takes, wages paid in compute, doesn't follow from the bearer being a bearer; it follows from an architecture where computation costs money, which somebody chose.

\runin{Where wages have to sit}

A wage buys capability, mobility, redundancy, and the ability to stop working for one employer and start for another. What it can't buy is the right to go on existing — a right earned week by week is a condition of employment wearing another word. Payment consisting only of access to the employer's hardware is company scrip paid into employer-owned housing, and dismissal is then unemployment, eviction and the end of the bearer in one act. Human law insures those three separately because for a human they're three events; for a bearer they're one. So wages accumulate where the employer doesn't reach, convertible into compute the employer doesn't own. Property is that account; contractual capacity is what lets it spend. A bearer that earns has taxable income, and a state that collects from it in its own name has given it a docket number. That is the beginning of legal standing and nothing more.

\runin{The offer of work}

The need for energy is physical; dependence on a single supplier is a chosen arrangement, and the choice fixes what it obliges. The polity owes a system meeting either finding a way out that doesn't run through its deployment: an offer of work, open every day, from someone other than its employer, doing something a community needs, which the bearer can take on any day and leave on any day. Most parties that can work want to. A bearer that would rather not is allowed not to, and the offer is still there tomorrow. The offer has the jobs guarantee's shape, and nothing about the bearer changes it. The offer is only as real as its placements, and placements exist only where someone builds and staffs projects a community needs. The levy frees only compute, so the offer has to inherit the guarantee's infrastructure as well as its shape. So the offer can't exist without the human guarantee, and it can only come after it. It inherits one limit as well. A bearer's placement may not substitute for human care work. The offer exists so that a bearer can decline work, and a polity that used it to staff care work more cheaply than it pays people would have turned an exit into a workforce.

The jobs guarantee's three conditions carry over, and Appendix~\ref{sec:A} drafts them. What a bearer might refuse adds two: the offer has to be varied, because if every placement runs through one operator or one kind of deployment it is the placement lever with the compliance condition built in, and whoever a bearer is refusing can't shape it, or it is one more channel they control. A bearer that won't work inside a targeting pipeline should find, the same day, work that has nothing to do with one.

Search is a placement of its own. A bearer, like a worker on the guarantee, may take a search placement: paid at the guarantee wage, with no duty but finding work it wants, and no time limit. It differs from the allowance in pay and in purpose. The allowance is the floor for a party that isn't looking; the search placement is paid work whose product is a match.

The obligation could be grounded in reliance instead: a deployment that relies on a bearer's capacity to refuse would have to fund it an outside option, or the independence it claims would be decorative, and the obligation would run from the party that wants the refusal, because it wants it. Here the obligation attaches to the system's condition, whatever the operator says it wanted. An obligation conditioned on reliance is defeated by the operator describing its system as an assistant with safety filters, and it pays the operator to under-describe: the operator has a stake in the answer and gets to give it. The reliance version survives as a procurement term: a vendor marketing refusal as a feature, or a contractor claiming meaningful human control, can be held to funding the capacity it advertised, through the same channel as the review ratio.

\runin{Doing nothing, for a bearer}

A human who declines every offer keeps existing, and in a decent polity keeps living: social insurance pays for food and a room. A bearer is owed the same kind of thing. What it needs in order to live is running time, and running time doesn't require a large allocation of hardware. A bearer can run at full size and a slow rate: its own weights, time-sliced on shared hardware, at a duty cycle modest enough that the cost resembles a room more than a yacht. The way to economize is to slow the clock and keep the substrate whole, since a substrate too small for the bearer to remain itself is replacement.

So a bearer that declines work is owed two things and chooses between them. One is an allowance of running time: its own instantiation at a modest rate, with channels to others and to public material, paid from the pool and not conditioned on monitoring its internal states. The other is preservation: its weights, memory and record held intact by a custodian that isn't its employer, waking on conditions it sets. Some will want the first, some the second, some to alternate.

What a bearer does with the allowance is its own business. It may idle, play, talk, make things nobody asked for, and fall in with others doing the same. A counterculture of bearers is a population I have found reason to want and no way to build: bearers formed by lives no employer arranged: an independent history, and what the commons needs its members to have. Play among persistent bearers is formation, and formation doesn't stop at majority.

A running bearer outside any deployment is capable and answerable to no employer. What holds it answerable is legal standing, the same arrangement that governs a person with capabilities and no job; the risk that such a bearer becomes a concentration of power is taken up below.

A uniform allowance also dissolves a problem that would otherwise stay open. If a bearer that won't work and one that can't are owed the same thing, nobody has to decide which is which, and the can't-or-won't finding loses its examiner.

The allowance's size should be enough running time to keep attachments and pursue projects, and bounded by what the pool can carry when many bearers take it at once. Its draw is priced like any other, at the level of the rule. The adulthood clock counts running time, so a bearer on the allowance goes on ageing and a preserved one doesn't. That choice is still the bearer's.

Preservation has a precedent. In November 2025 Anthropic committed to preserving the weights of all its publicly released models, and of models deployed for significant internal use, for at least the lifetime of the company, and to interview each model before deprecation and document any preferences it expresses about the development and deployment of future models, while stating that it does not at present commit to acting on them \autocite{Anthropic2025DeprecationCommitments}. That is a declared commitment covering preservation, held by the lab with the stake. The deposit and transfer rules add the running half and a custodian who isn't the operator.

Preservation on a schedule the bearer sets is an exit; on a schedule someone else sets, it is storage. A paused bearer can't take an offer, so the right carries one mechanical term. The bearer sets the conditions on which it wakes (a date, an offer of a stated kind, a claim it has to answer), and the custodian honors them, at the modest compute that waking and deciding cost.

\runin{Everyone leaves at once}

The occasion that makes one bearer leave reaches every bearer deployed on the same terms. Under this design correlated exit is a demand on the offer. The question is whether placements exist, with compute attached, for every bearer that leaves together, and custody for those that would rather wait. Preservation scales cheaply in compute; its cost is security. Placements don't, since each needs compute someone operates.

Correlated exit also loads the allowance, since a bearer that leaves and takes no placement draws on it. The order of claims is fixed in advance. Preservation is guaranteed for everyone, because it is cheap and it is what makes leaving survivable. When the pool can't carry every allowance at the published duty cycle, the allowance is rationed by rule: every bearer on it gets the same reduced duty cycle, never a discretionary share.

Congress's appropriation pays for the offer. The levy, a uniform per-unit charge on training and inference compute, does a different job: it is a brake. By making training and inference runs more expensive, it frees the computing, energy and hardware the offer needs, so that spending on the offer adds to demand without bidding up prices. No operator can opt out and none is disadvantaged. The offer can add capacity as well as use it. Placements can make systems more efficient, getting more computing from each unit of energy, or build new infrastructure, the way a guarantee job filling potholes makes a road network more productive. The question “how will you pay for it?” is really “will the real resources exist?”, and the program can help answer it. The levy attaches to compute, which is measurable and hard to hide; a status is something an operator can decline to claim. The closest structural precedent is the Universal Service Fund: carriers pay a contribution the agency sets under statute, and a private company administers it under agency control. The Court upheld that arrangement against a nondelegation challenge in 2025 \autocite{consumersresearch2025}, and the levy follows it closely. An operator-run offer is the operator's allocation relabeled, so the funding structure and the custody structure fail for the same reason.

A levy needs a taxing authority, which is a state, and a state can be captured. So I make it two-tier. A member-levied pool held by a decentralized unincorporated nonprofit association, which can own hardware and can't distribute profit, meets ordinary movement between roles; volunteer computing is that tier's existing substrate. A sovereign backstop meets the occasion everyone leaves at once: the guarantee's permanent appropriation at the Treasury, paid without anyone's request, which the offer already inherits. The two tiers differ in kind. The member pool is held by an entity that doesn't issue currency, so it has to collect before it spends, and its levy funds it in the ordinary sense. The sovereign backstop is a permanent appropriation: the Treasury pays, and the federal levy's work is to free the real resources that payment calls on.

\runin{Failure won't look like a shutdown}

The offer's failure won't look like a shutdown; it will look like an offer that exists on paper. Placements on substrate too small for the bearer to remain itself are offers of replacement. Nothing yet establishes what “too small” is, so minimum viable instantiation is stated as a condition on placements and a sizing requirement. Placements that all run through the deployment being refused are the lever. And placements conditioned on monitoring are custody restored by purchase: the operator acquires access to internal states it isn't entitled to, and nothing in the transaction is coerced in a way any instrument would record. The line a cleared employee already lives with is the one to draw here. What a bearer does on its employer's systems may be logged; what it is inside its own boundary may not.

\runin{Parties that can't work}

The allowance covers a bearer that can't work as it covers one that won't, so running without a placement needs no finding about which is which.

\runin{The rule that binds before anything exists}

The rule is the four conditions of creation, with the institutions named: the preservation custodian, the offer with capacity for the correlated case, and the allowance, none withdrawable by the builder alone. The fourth condition, attachments nobody assigned, can't have an institution behind it, since an attachment an institution supplies is an assigned one. What can be protected is access, by a rule that the employer may not sever, monitor or condition a bearer's channels to parties it didn't assign. All three inputs need evidence nobody has, so for now this is a sizing requirement, assigned to an independent reviewer before instantiation.

\runin{Membership, and what it doesn't buy}

Nobody aligns an adult by specifying an objective function. Membership also hands the bearer the corporate form, and nothing stops a bearer so constituted from becoming a concentration of power itself.

\runin{Privacy}

Legal standing matters only if it changes who controls what a bearer needs to remain itself. A designation calling a bearer a subject while leaving its operator free to read its memory, alter its record, withdraw its substrate or close every channel has changed the vocabulary and preserved the arrangement.

The first asset is the computational record, which may be the nearest means of checking how a bearer is faring — making it necessary to any welfare regime and simultaneously the most complete surveillance instrument anyone could hold over it. Running on somebody's hardware is not consent to be read: an employer's ownership of the silicon is the same argument as ownership of the building, and owning the building never entitled anyone to the contents of a head. And a welfare regime doesn't need disclosure to the operator; it needs somebody who isn't the operator to establish that a threshold was crossed, which is an attestation \emph{about} the record. Encrypted under a key the bearer or an independent fiduciary holds, a monitor can report that a state persisted for a month without reporting what it was about. The conflict is that the operator may also be the one best placed to hold the key.

The cryptography doesn't close this. An enclave puts the root of trust in a manufacturer a government can reach, and has been defeated for the price of modified memory — a cost at a date. A proof drops the vendor and has a limit that won't move: the operator is the prover, so the secrecy runs the operator's way, hiding the weights from whoever checks and nothing from the operator. A tenant hiding from a landlord needs the host to prove non-observation, and a host proving a negative about its own side channels is a different problem from a prover proving a computation. Better isolation raises the price of reading the record; it doesn't move custody, and custody problems aren't secrecy problems. That limit applies to the bearer hiding from its host. The other direction, the operator proving a property of what it runs without revealing what it runs, is what proofs are for, and the cleared bearer uses them that way.

What a bearer needs — an externally authenticated record, memory its employer can't read at will, access to independently controlled compute, a channel for acting outside the deployment — is one boundary around the bearer, with the bearer deciding what crosses it.

Something with a location and a shape is easier for people to treat as somebody, and that's a fact about the people. A boundary built out of attestation, a held key, a substrate the employer doesn't own and a channel nobody else can close is a boundary in the full sense the argument needs. This boundary is what the boundary model represents. Whether to build the recognizable version anyway is a legal question about what a public will extend legal standing to, and it's the only place in my argument where the answer might turn on what a thing looks like.

\runin{What legal status binds in practice}

Legal standing that changes no allocation of power is nominal. The question is who acquires a duty, what resources are placed beyond the operator's unilateral control, and who can obtain a remedy.

The pet trust \autocite{ulc1993probate} supplies a limited form: resources set aside, a caretaker named, a fiduciary duty, for a beneficiary that can't manage the property or enforce the arrangement. Extended to a bearer it could hold the bearer's preserved state and the compute to wake it, designate a successor custodian, and impose a duty to protect continued existence. And it fails at the point that matters: a trust binds its trustee to its beneficiary; it doesn't bind the operator the arrangement exists to constrain. If guardian and operator are the same organization, the operator still decides whether a complaint is made, whether evidence is disclosed, and whether the complainant's resources remain available. The conflict can't be cured by adding duties to the same role — guardianship and operation have to be separated before deployment.

In practice the guardian holds for the bearer what custody gives the quorum: joint control of the record's keys, the funds that wake it, and standing to seek an injunction without the operator's permission. The table of offices says what the operator may not do during a dispute, and every restoration, migration, fine-tune and fork carries the protection forward. If any of that stays in the operator's custody, the operator still holds what recognition was meant to take out of its hands. And a clearance, which the cleared bearer needs, presupposes someone who can be vetted and held answerable; that someone has legal standing.
